\documentclass[../../../include/open-logic-section]{subfiles}

\begin{document}

\olfileid{sth}{ordinals}{vn}	
\olsection[د فون نويمن جوړونه]{د فون نويمن له خوا د ترتيبي عددونو جوړونه}

\olref[sth][ordinals][iso]{thm:woalwayscomparable} يو فکر راولاړوي.
موږ ځينې شيان معرفي کولے شو چې \emph{د ترتيب ډولونه} ئې وبولو او د
ښه ترتيبونو استازيتوب پرې وکړو. که $\ordtype{A, <}$ د ښه ترتيب $\tuple{A, <}$ د ترتيب ډول
په توګه وليکو، نو غواړو دا دوه اصلونه ترلاسه کړو:
\begin{align*}
	\ordtype{A, <} = \ordtype{B, \lessdot} &
	\text{ هله او يوازې هله چې } \ordeq{\tuple{A, <}}{\tuple{B, \lessdot}}\\
	\ordtype{A, <} < \ordtype{B, \lessdot}&
	\text{ هله او يوازې هله چې }\ordeq{\tuple{A, <}}{\tuple{B_b, \lessdot_b}}\text{ د يو مناسب }b \in B
\end{align*}
همدارنګه، هيله لرلے شو چې د ترتيب ډولونه \emph{د ځينو ټاکلو سټونو په توګه}
معرفي کړو، لکه څنګه چې طبيعي عددونه د ځينو ټاکلو سټونو په توګه معرفي کولے شو.

د دې تر ټولو عامه طريقه، او هغه لاره چې موږ به ئې تعقيب کړو، دا ده چې
د ترتيب دا ډولونه د ځينو \emph{کانوني} ښه مرتب شويو سټونو له لارې تعريف کړو.
دا کانوني سټونه په لومړي ځل فون نويمن معرفي کړي وو:

\begin{defn}
سټ $A$ \emph{متعدي} دے، {که او يوازې که} $(\forall x \in A)x \subseteq
A$ وي. بيا $A$ يو \emph{ترتيبي عدد} دے، {که او يوازې که} $A$ متعدي وي او
$\in$ ئې ښه ترتيب ورکړي.
\end{defn}
\noindent
په لاندې بحث کښې به د ترتيبي عددونو دپاره يوناني توري کاروو. له تعريف څخه
سمدستي نتيجه اخلو چې، که $\alpha$ ترتيبي عدد وي، نو
$\tuple{\alpha, \in_\alpha}$ يو ښه ترتيب دے، چې پکښې $\in_\alpha =
\Setabs{\tuple{x, y} \in \alpha^2}{x \in y}$ دے. نو د نښانو په کارولو کښې
په لږې بې احتياطۍ سره ويلے شو چې $\alpha$ \emph{پخپله} ښه ترتيب دے.

زموږ لومړني څو ترتيبي عددونه دا دي:
\[
	\emptyset, \{\emptyset\},
	\{\emptyset, \{\emptyset\}\},
	\{\emptyset, \{\emptyset\}, \{\emptyset, \{\emptyset\}\}\}, \ldots
\]
و به وينئ چې دا هماغه لومړني څو ترتيبي عددونه دي چې زموږ د نامتناهيت په
بديهي اصل کښې مو ليدلي وو، يعنې د فون نويمن د $\omega$ په تعريف کښې
(\olref[sth][z][infinity-again]{sec} وګورئ). دا تصادف نۀ دے.
د فون نويمن د ترتيبي عددونو تعريف طبيعي عددونه ترتيبي عددونه ګڼي، خو
د نامتناهيت هاخوا ترتيبي عددونو ته هم ځای ورکوي.

د تل په شان اوس پوښتلے شو: ايا دا \emph{هماغه} ترتيبي عددونه دي؟
يا فون نويمن يوازې ځينې داسې سټونه راکړي دي چې د ترتيبي عددونو په توګه
ئې \emph{کارولے} شو؟ د دې پوښتنې بحثونه هغو بحثونو ته ورته دي چې په
\olref[sfr][rel][ref]{sec}،
\olref[sfr][arith][ref]{sec}،
\olref[sfr][infinite][dedekindsproof]{sec} او
\olref[sth][z][nat]{sec} کښې مو کړي وو؛ نو خبره به نۀ اوږدوو.
پر ځای ئې، په لاندې بحث کښې به «ترتيبي عددونه» د «فون نويمن ترتيبي
عددونو» دپاره کاروو.

\end{document}
